\documentclass[10pt]{article}
\usepackage[T1]{fontenc}
\usepackage[utf8]{inputenc}
\usepackage[french]{babel}
\usepackage{amsmath,amssymb,mathrsfs,mathpazo}
\usepackage{geometry}
\geometry{a4paper,margin=2cm}
\usepackage{graphicx,array,tabularx,tabularray,colortbl,multirow,multicol,booktabs,enumitem,pifont,ulem,fancybox,diagbox,cancel}
\usepackage{pstricks,pstricks-add,pst-plot,pst-tree,pst-node,pst-text,pst-func,pst-math,pst-bspline,pst-3dplot,pst-eucl,pst-solides3d}
\let\pscalculate\undefined
\usepackage{pst-bezier}
\usepackage{tikz,pgfplots,tkz-tab}
\usetikzlibrary{arrows,arrows.meta,patterns,calc,positioning,trees,decorations.markings,shapes.geometric,shapes.misc,shapes.arrows,intersections,decorations.pathmorphing,decorations.pathreplacing,backgrounds,shadows,babel}
\pgfplotsset{compat=1.18}
\usepgfplotslibrary{fillbetween}
\usepackage[np]{numprint}
\usepackage{listings,makecell,pas-tableur,twoopt}
\let\barre\undefined
\usepackage[pstricks]{tablvar}
\definecolor{darkgreen}{RGB}{0,100,0}
\providecommand{\red}{\color{red}}
\providecommand{\blue}{\color{blue}}
\providecommand{\green}{\color{green}}
\providecommand{\black}{\color{black}}
\providecommand{\white}{\color{white}}
\providecommand{\gray}{\color{gray}}
\providecommand{\cyan}{\color{cyan}}
\providecommand{\magenta}{\color{magenta}}
\providecommand{\yellow}{\color{yellow}}
\providecommand{\violet}{\color{violet}}
\providecommand{\orange}{\color{orange}}
\providecommand{\euro}{\texteuro}
\providecommand{\EUR}{\texteuro}
\providecommand{\ieme}{\textsuperscript{e}}
\providecommand{\degres}{\ensuremath{^\circ}}
\providecommand{\degree}{\ensuremath{^\circ}}
\providecommand{\vv}[1]{\overrightarrow{#1}}
\providecommand{\Coord}[1]{#1}
\providecommand{\decosix}{}
\providecommand{\starredbullet}{\ensuremath{\bullet}}
\pagestyle{empty}
\setlength{\parindent}{0pt}
\usepackage{adjustbox}\DeclareUnicodeCharacter{2236}{:}\DeclareUnicodeCharacter{2190}{\ensuremath{\leftarrow}}
\providecommand{\R}{\mathbb{R}}
\providecommand{\N}{\mathbb{N}}
\providecommand{\Z}{\mathbb{Z}}
\providecommand{\Q}{\mathbb{Q}}
\providecommand{\D}{\mathbb{D}}
\providecommand{\C}{\mathbb{C}}
\providecommand{\e}{{\mathrm{e}}}
\providecommand{\vect}[1]{\overrightarrow{#1}}
\providecommand{\vectt}[1]{\overrightarrow{\mathrm{#1}}}
\providecommand{\vvt}[1]{\overrightarrow{\mathrm{#1}}}
\providecommand{\np}[2][]{{#2~#1}}
\providecommand{\vv}[1]{{\overrightarrow{#1}}}
\renewcommand{\d}{{\mathrm{d}}}
\providecommand{\degres}{{^\circ}}
\providecommand{\degre}{{^\circ}}
\providecommand{\blue}{\color{blue}}
\providecommand{\red}{\color{red}}
\providecommand{\enskip}{\quad}
\providecommand{\llbracket}{\lbrack\!\lbrack}
\providecommand{\rrbracket}{\rbrack\!\rbrack}
\providecommand{\euro}{{\text{€}}}
\providecommand{\no}{{\text{n°}}}
\providecommand{\interval}[2]{{[#1;#2]}}
\providecommand{\intervFO}[2]{{[#1;#2[}}
\providecommand{\dotfill}{\ldots}

\providecommand{\ds}{\displaystyle}
\providecommand{\textding}[1]{#1}
\providecommand{\Oij}{\(\mathrm{O};\vec{\imath},\vec{\jmath}\)}
\providecommand{\Oijk}{\(\mathrm{O};\vec{\imath},\vec{\jmath},\vec{k}\)}
\providecommand{\Ouv}{\(\mathrm{O};\vec{u},\vec{v}\)}
\providecommand{\intervFF}[2]{\ensuremath{[#1;#2]}}
\renewcommand{\euro}{\ifmmode\text{\texteuro}\else\texteuro\fi}
\DeclareUnicodeCharacter{00B0}{\ifmmode^{\circ}\else\textdegree\fi}
\DeclareUnicodeCharacter{20AC}{\ifmmode\text{\texteuro}\else\texteuro\fi}
% Énoncés : Éducation nationale. Transcription et corrigés : APMEP, Denis Vergès et contributeurs. Découpage et classement : CoopMaths, Sébastien Lozano, Éric Elter et contributeurs. Adaptation de la présentation : Maths Explorer.
% CC BY-SA : https://coopmaths.fr/www/about/
% https://coopmaths.fr/alea/static/e3c/2020/tex/e3c_2020_00_sujet20_3.tex
\begin{document}
\noindent\textit{APMEP / CoopMaths. CC BY-SA. Présentation adaptée par Maths Explorer.}\par\medskip
\medskip

La famille A décide de diminuer de 2\,\% par mois sa quantité de déchets produits par mois à partir du 1\up{er} janvier 2020.

Au mois de décembre 2019, elle a produit $120$~kg de déchets.

\medskip

\begin{enumerate}
\item Justifier qu'au bout de 2 mois, la famille A aura produit environ $115$~kg de déchets.
\end{enumerate}

On admet que la quantité de déchets produits chaque mois conserve la même évolution toute l'année.

On modélise l'évolution de la production de déchets de la famille A par la suite de terme général $a_n$, où $a_n$ représente la quantité, en kg, de déchets produits par la famille A $n$ mois après décembre 2019.

Ainsi, $a_0$ représente la quantité de déchets produits durant le mois de décembre 2019, $a_1$ représente la quantité de déchets produits durant le mois de janvier 2020, etc.

\begin{enumerate}[resume]
\item  
	\begin{enumerate}
		\item Déterminer la nature de la suite $\left(a_n\right)$.
		\item Pour tout entier naturel $n$, exprimer $a_n$ en fonction de $n$.
		\item Déterminer la quantité totale de déchets que produira la famille A durant l'année
2020.

On arrondira le résultat à l'unité.

On rappelle que :

Soit $\left(a_n\right)_{n\in \N}$ une suite géométrique de raison $q$,\,$ q \ne 1$. La somme $S$ de termes consécutifs est égale à $S = u_1 + u_2 +\ldots + u_n = u_1 \times \dfrac{1 - q^n}{1 - q}$.
		\item On donne le programme ci-dessous.
		












 
\begin{center}
\begin{tabular}{c l}
\texttt 1 & \texttt{{\blue def} S(n) :}\\
\texttt 2 &  \vline \quad \texttt{U = 120}\\
\texttt 3 &  \vline \quad \texttt{S = 0}\\
\texttt 4 &  \vline \quad \texttt{{\blue for} k {\blue in} range (n) :}\\
\texttt 5 &  \vline \quad \vline \quad \texttt{U = 0.98 * U}\\
\texttt 6 &  \vline \quad \vline \quad \texttt{S = S + U}\\
\texttt 7 &  \vline \quad \texttt{{\blue return} (S)}\\
\texttt 8 & 
\end{tabular}
\end{center} 
 
Que représente le résultat renvoyé par la fonction si on entre l'instruction S(6) ?
	\end{enumerate}
\end{enumerate}

\bigskip

\end{document}
